\section{Embedded Software Architecture}

The firmware running on the STM32 microcontroller was developed using a modular, object-oriented approach in
C++. This architecture ensures code reusability, encapsulation of hardware specifics, and ease of maintenance
and scalability. The software stack (Appendix \ref{app:software}) is organized into four distinct abstraction
layers, moving from low-level hardware management to high-level application logic.


\subsubsection{Layer 1: Vendor HAL}

The foundational layer consists of the Hardware Abstraction Layer (HAL) provided by the microcontroller
manufacturer. Modules such as \texttt{i2c.h} and \texttt{tim.h} handle the direct register-level configuration
of the MCU peripherals. Relying on these standard verified drivers ensures robust handling of interrupts and
basic I/O operations.


\subsubsection{Layer 2: Custom Middleware and Communication}

This layer acts as a bridge between the raw hardware and the high-level logic, abstracting the communication
protocols. Key modules include:

\begin{itemize}
    \item \textbf{I2C Device Wrappers:} Classes like \texttt{i2c\_device} and \texttt{i2c\_slave} extend the
        Vendor HAL to provide simplified interfaces for bus communication, handling read/write buffers and
        addressing logic.
    \item \textbf{Virtual Register:} This is the core component of the communication architecture. It
        implements the "Virtual Register Map" concept, abstracting the internal state of the actuator
        (positions, gains, modes) as a set of memory-mapped values. This allows the external master to
        interact with the servo seamlessly, as if reading from a standard memory device, changing its
        behaviour dynamically as needed.
\end{itemize}


\subsubsection{Layer 3: High-Level Component Drivers}

This layer contains the specific drivers for the chosen hardware components and the pure algorithmic logic,
independent of the specific MCU.

\begin{itemize}
    \item \textbf{Component Drivers:} The \texttt{AS5600} and \texttt{DRV8870} classes encapsulate the
        specific operational details of the sensor and driver. They utilize the lower layers to perform tasks
        such as reading raw magnetic angles or setting PWM duty cycles, exposing clean methods (e.g.
        \texttt{getAngle()}, \texttt{setSpeed()}) to the application layer.
    \item \textbf{Math and Control:} This module contains the implementation of the control algorithms,
        including the Switching PID logic. It is designed to be purely mathematical, accepting error inputs
        and producing control outputs without direct dependencies on hardware headers.
\end{itemize}


\subsubsection{Layer 4: Application Layer}

The highest level of the hierarchy integrates all previous modules to define the actuator's behavior.

\begin{itemize}
    \item \textbf{Smart Servo Class:} This main class encapsulates the entire logic of the actuator. It
        instantiates the drivers and control modules, managing the data flow between the sensor readings,
        the PID computation, and the motor output.
    \item \textbf{Entry Point:} The \texttt{main.cpp} and \texttt{CONFIG.h} files handle the system
        initialization, configuration, and the execution of the main real-time control loop.
\end{itemize}


\subsection{Control Strategy}

Achieving high-precision positioning in robotic joints is inherently challenging due to the presence of
nonlinear friction phenomena \cite{armstrong1994survey}, specifically static friction (Coulomb) and velocity
weakening effects (Stribeck). As extensively discussed in control literature, standard linear controllers like
the classical PID struggle to manage these non-linearities effectively.

In particular, the integral action required to overcome static friction often leads to overshoot. When
combined with the Stribeck effect, this energy imbalance prevents convergence to the setpoint, resulting in
persistent oscillations known as the "hunting phenomenon". While model-based control techniques could
theoretically compensate for these effects, they require precise parameter identification — which is difficult
to maintain over time due to wear and load changes — and impose a computational burden often incompatible with
low-cost embedded hardware.

A robust alternative that seems to mitigate these issues without requiring complex friction models is the
class of hybrid control strategies, such as Reset or Switching PID \cite{bisoffi2020stick, lorigiola2025switching}.
These approaches modify the internal state of the controller (specifically the integrator) based on discrete
logical rules triggered by the system's state (e.g. zero-crossing of the error or velocity).

The control strategy implemented in this custom actuator is directly derived from the Switching PID strategy
proposed and analyzed in \cite{lorigiola2025switching}. This approach uses a state-machine logic to
dynamically switch the integrator behavior and avoid hunting.

Since the rigorous mathematical derivation and stability proofs of this control strategy are extensive and already
well-documented in the referenced literature, they are considered outside the scope of this thesis. Instead,
this section provides a functional description of the algorithm and the pseudocode implementation specifically
tailored for the custom hardware application.


\subsubsection{Algorithm Description}

\begin{algorithm}[H]
    \footnotesize
    \caption{Switching PID for Friction Compensation}
    \label{alg:switching_pid}
    \begin{algorithmic}[1]
        \Require Reference $r_k$, Measurement $y_k$, Gains $K_p, K_i, K_d$
        \Ensure Control Input $u_k \in [-1, 1]$

        \State $e_k \gets r_k - y_k$ \Comment{Position Error}
        \State $v_k \gets \frac{y_k - y_{k-1}}{\Delta t}$ \Comment{Estimated Velocity}

        \vspace{0.1cm}

        % --- BLOCCO 1: PI Computation (Blu) ---
        \ColorBlock{RoyalBlue}{
            \textbf{1. PI Computation \& Anti-Windup} \\
            $\sigma \gets K_i \cdot e_k$ \\
            $\phi' \gets K_p \cdot (e_k + I_{state})$ \\
            $\phi \gets \text{Saturation of }\phi'$ \\
            $\Delta_{sat} \gets \phi' - \phi$ \\
            \textbf{if} $\text{sgn}(\sigma) == \text{sgn}(\Delta_{sat})$ \textbf{and} $\Delta_{sat} \ne 0$ \textbf{then} \\
            \hspace*{1.5em} $I_{state} \gets I_{state}$ \Comment{Hold integrator} \\
            \textbf{else} \\
            \hspace*{1.5em} $I_{state} \gets I_{state} + \sigma \Delta t$ \Comment{Update integrator} \\
            \textbf{end if}
        }

        \vspace{0.1cm}

        % --- BLOCCO 2: Switching Logic (Arancione) ---
        \ColorBlock{OrangeRed}{
            \textbf{2. Switching Logic} \\
            \textbf{if} $(\text{State} == \textsc{StickSlip}) \textbf{ and } (\sigma \cdot \phi \le 0)$ \textbf{then} \\
            \hspace*{1.5em} $\phi \gets -\phi$ \Comment{Reset energy} \\
            \hspace*{1.5em} $I_{state} \gets \frac{\phi}{K_p} - e_k$ \\
            \hspace*{1.5em} $\text{State} \gets \textsc{Overshoot}$ \\
            \textbf{else if} $(\text{State} == \textsc{Overshoot}) \textbf{ and } (v_k \cdot \phi \ge 0)$ \textbf{then} \\
            \hspace*{1.5em} $\phi \gets K_p \cdot \sigma / K_i$ \Comment{Restore linear behavior} \\
            \hspace*{1.5em} $I_{state} \gets 0$ \\
            \hspace*{1.5em} $\text{State} \gets \textsc{StickSlip}$ \\
            \textbf{end if}
        }

        \vspace{0.1cm}

        % --- BLOCCO 3: Output (Verde) ---
        \ColorBlock{ForestGreen}{
            \textbf{3. Output Calculation} \\
            $u_k \gets \text{Clamp}(\phi - K_d \cdot v_k, -1, 1)$
        }

        \vspace{0.1cm}

        \State \Return $u_k$
    \end{algorithmic}
\end{algorithm}

The control loop runs at a fixed frequency of $1\,\text{kHz}$. At each time step $k$, the algorithm described
in Algorithm \ref{alg:switching_pid} is executed.

The controller operates in two distinct logical states: when in \textbf{Stick-Slip State} (default) the
integrator is fully active to build up the torque necessary to overcome static friction ($F_s$). The logic
prioritizes "breaking out" of the stick phase. Detecting that the error has crossed zero (overshoot) triggers
the transition to the \textbf{Overshoot State}. In this mode, the logic serves to dampen the response by
inverting the control action, preventing the limit cycles (hunting) typical of high-gain PIDs on frictional loads.

The transition between these states is governed by the relationship between the generalized position error
($\sigma$) and the total proportional-integral control effort ($\phi$).

Notably, the term $\phi$ acts as a 'virtual' control effort that is reset discontinuously during state
transitions (Block 2), allowing the controller to instantly adapt its energy based on the
friction regime, a capability not present in standard linear PIDs.

This control strategy is originally derived for set-point regulation (point-to-point motion). While effective for
the quasi-static segments of the gait cycle, its application to dynamic trajectory tracking may require further
tuning or structural modifications (e.g. feed-forward terms). Addressing these specific dynamic tracking
optimization is considered outside the scope of this thesis and is left for future work.



